\chapter{Testing and Results}

Everything in this chapter is measured on the held-out 20\% test set --- the same 102 unseen contracts for every model. The numbers were regenerated by a standalone evaluation notebook that loads the trained models from disk, so no state left over from the training session can influence them.

\section{Presence Classification}

\subsection{Model comparison}
Table~\ref{tab:ensemble} and Figure~\ref{fig:ensemble} compare all the presence models on the 4{,}182 (contract, category) test decisions. Two things stand out. First, the windowed transformer \emph{loses on its own} to the full-document TF--IDF baseline (0.694 vs.\ 0.775). Legal vocabulary is distinctive enough that lexical evidence over the whole document beats semantic evidence over windows, and max-pooling inflates false positives exactly the way Section~\ref{sec:mil} predicted. Second, the AND-ensemble --- which reports a clause only when \emph{both} models agree --- filters out each model's uncorrelated false positives and finishes on top (0.776). The margin over the baseline is small on 102 test contracts, so we do not lean on it: Section~\ref{sec:bootstrap} quantifies exactly how small, and finds that this particular gap is indistinguishable from zero. The conclusion we stand behind is methodological: \emph{simple lexical models are strong at document-level presence, and the transformer earns its place as a precision filter, not as a replacement}.

\begin{table}[ht]
\centering
\caption{Presence models on the 20\% test set.}
\label{tab:ensemble}
\begin{tabular}{lc}
\toprule
\textbf{Model} & \textbf{micro-F1}\\
\midrule
\textbf{Ensemble AND (both must agree)} & \textbf{0.776}\\
TF--IDF baseline (full document)        & 0.775\\
Ensemble AVG (mean probability)         & 0.718\\
Ensemble OR (either fires)              & 0.696\\
Transformer (max-pool) alone            & 0.694\\
\addlinespace
\multicolumn{2}{l}{\emph{control run, Section~\ref{sec:budget}:}}\\
Transformer trained on all 53{,}662 windows & 0.694\\
Ensemble AND on that transformer            & 0.782\\
\bottomrule
\end{tabular}

{\footnotesize\singlespacing \emph{Note.} micro-F1 over all 4{,}182 test cells.\par}
\end{table}

\begin{figure}[ht]
\centering
\includegraphics[width=0.75\textwidth]{ensemble_comparison.png}
\caption{micro-F1 of all presence models on the test set.}
\label{fig:ensemble}
\end{figure}

\subsection{How much of that ranking is real?}
\label{sec:bootstrap}
Table~\ref{tab:ensemble} is a leaderboard, and leaderboards invite a reader to treat the ordering itself as a result. On 102 test contracts a margin of 0.776 against 0.775 --- roughly one decision in a thousand --- plainly cannot support that reading, so rather than leave the caveat as prose we measured it.

The natural tool is a \textbf{bootstrap confidence interval}, but the resampling unit matters. The 4{,}182 test decisions are not independent: 41 of them come from the same contract, and a contract that is unusual --- unusually long, unusually drafted --- moves all 41 together. Resampling individual (contract, category) cells would ignore that structure and produce intervals that are far too narrow. We therefore resample \emph{contracts} with replacement (a cluster bootstrap), drawing 102 contracts from the 102 available, recomputing micro-F1 for both models on each resample, and repeating 10{,}000 times. Both models are scored on exactly the same resample every time, so the difference is \emph{paired} --- which is what allows a difference interval much tighter than the two individual intervals would suggest.

\begin{table}[ht]
\centering
\caption{Cluster bootstrap over the 102 test contracts.}
\label{tab:bootstrap}
\small
\begin{tabular}{lccc}
\toprule
\textbf{Comparison} & \textbf{Difference in micro-F1} & \textbf{95\% CI} & \textbf{$P(>0)$}\\
\midrule
AND-ensemble vs.\ TF--IDF baseline & $+0.0014$ & $[-0.0068,\ +0.0089]$ & 0.626\\
TF--IDF baseline vs.\ transformer  & $+0.0805$ & $[+0.0626,\ +0.0992]$ & 1.000\\
AND-ensemble vs.\ transformer      & $+0.0818$ & $[+0.0654,\ +0.0987]$ & 1.000\\
\bottomrule
\end{tabular}

{\footnotesize\singlespacing \emph{Note.} 10{,}000 resamples. The paired difference is what settles each comparison; $P(>0)$ is the fraction of resamples in which the first model wins.\par}
\end{table}

The three comparisons do not come out the same way, and the difference between them is the point.

\textbf{The ensemble's win over the baseline is not real.} The interval is $[-0.0068, +0.0089]$: it straddles zero, and the ensemble is ahead in only 62.6\% of resamples, barely better than a coin flip. Taken individually the two models score 0.776 (95\% CI $[0.759, 0.792]$) and 0.775 (95\% CI $[0.758, 0.791]$), intervals that almost entirely overlap. On this test set the AND-ensemble and the TF--IDF baseline are \emph{statistically indistinguishable}, and we now state that as a measurement rather than as a hedge. The top two rows of Table~\ref{tab:ensemble} should be read as a tie.

\textbf{The transformer's loss, by contrast, is real.} The baseline beats the windowed transformer by 0.0805, with an interval of $[+0.0626, +0.0992]$ that is nowhere near zero, and it wins in every one of the 10{,}000 resamples. This matters for how the thesis' two empirical claims should be weighted, because they are not on equal footing. The finding we stand behind most firmly --- that a full-document lexical model beats a windowed DistilBERT at document-level presence classification --- is the one that clears the significance test comfortably. The finding we were already cautious about --- that the conjunctive ensemble edges past the baseline --- is the one that does not. That our caution happened to point the right way is reassuring, but it was worth checking rather than assuming, and a reader should not have had to take it on trust.

None of this undercuts the case for the ensemble, because that case was never the 0.001 of micro-F1. The AND rule earns its place through the \emph{error profile} it produces rather than its position on the leaderboard: it converts the transformer's inflated false positives into balanced precision and recall, which is what the next section examines.

\subsection{How much of this moves if we only change the random seed?}
\label{sec:seeds}
Section~\ref{sec:bootstrap} measured one kind of uncertainty --- how much the score depends on \emph{which contracts} landed in the test set --- and used it to call the ensemble's margin a tie. There is a second kind it cannot see. Every number reported so far comes from a single training run per model, and a fine-tuning run is stochastic: weight initialization of the classification head, data shuffling, and the negative-window sample all depend on the seed. If simply reseeding moves the score as much as the effects we are discussing, then a leaderboard built from single runs is reporting noise regardless of how the test set was drawn.

This is easy to check and we had not checked it, so we retrained the presence transformer twice more at the same 24{,}000-example budget, changing nothing but the seed (42, 43, 44), and re-evaluated each on the same 102 contracts. Table~\ref{tab:seeds} reports the spread.

\begin{table}[ht]
\centering
\caption{The presence transformer across three random seeds.}
\label{tab:seeds}
\small
\begin{tabular}{lcccccc}
\toprule
\textbf{Model} & \textbf{Seed 42} & \textbf{Seed 43} & \textbf{Seed 44} & \textbf{Mean} & \textbf{SD} & \textbf{Range}\\
\midrule
Transformer (max-pool) & 0.6943 & 0.6928 & 0.7082 & 0.6984 & 0.0085 & 0.0154\\
AND ensemble           & 0.7761 & 0.7774 & 0.7803 & 0.7780 & 0.0022 & 0.0042\\
\midrule
TF--IDF baseline       & \multicolumn{3}{c}{0.7747 (deterministic)} & 0.7747 & --- & ---\\
\bottomrule
\end{tabular}

{\footnotesize\singlespacing \emph{Note.} All three runs use the 24{,}000-example budget and differ only in seed; the TF--IDF baseline is deterministic and appears once.\par}
\end{table}

\textbf{The ensemble's margin is smaller than its own seed noise.} The AND ensemble's advantage over the baseline, the number Table~\ref{tab:ensemble} ranks it by, is $+0.0014$. The standard deviation of that same ensemble across seeds is $0.0022$, and the observed range is $0.0042$ --- three times the margin. Reseeding the model moves the score further than the result being claimed for it. This is an entirely independent line of evidence pointing where the cluster bootstrap already pointed: had we trained with seed 44 and stopped, we would have reported a margin of $+0.0056$ and been just as wrong to make anything of it. A single-run leaderboard separated by one part in a thousand was never measuring what it appeared to measure.

\textbf{The finding we do stand behind is untouched by seed noise.} The baseline beats the transformer by 0.0805, which is 9.5 times the transformer's seed-level standard deviation of 0.0085. Even the luckiest of the three seeds, at 0.7082, still loses to the deterministic baseline by 0.0665. There is no seed at which the windowed transformer becomes competitive, and combining this with Section~\ref{sec:budget} --- where doubling the training data also failed to close the gap --- leaves the architectural explanation as the one standing after both alternatives were tested.

It is worth noticing which model is noisier, because it is not the one a reader might guess. The transformer alone varies four times as much across seeds ($\mathrm{SD} = 0.0085$) as the ensemble containing it ($0.0022$). Requiring agreement with a deterministic lexical model damps the transformer's run-to-run variation as well as its false positives: a spurious detection has to survive a second, fixed opinion, and seed-driven fluctuations in what the transformer fires on are largely filtered out. That stabilization is a real and previously unreported property of the conjunctive rule --- and, on the evidence of Section~\ref{sec:raretail}, one it buys at the price of rare-category recall. It is a better argument for the ensemble than its position on the leaderboard ever was.

\textbf{Scope of this check.} We reseeded the presence transformer only. The span extractor and the summarizer remain single runs, so their reported scores carry an unmeasured variance component of unknown size, and we flag that in Section~8.2 rather than assume it is small.

\subsection{Was the transformer simply undertrained?}
\label{sec:budget}
There is an objection to Section~\ref{sec:bootstrap} that the bootstrap cannot answer, and it goes to the heart of the thesis' central empirical claim. The TF--IDF baseline is fitted on all 408 training contracts in full. The window-level transformer was not: to keep the run tractable on a laptop it was trained on a stratified subsample of 24{,}000 of the 53{,}662 available window examples (Section~5.2). So when the transformer loses by 0.081 micro-F1, we cannot tell from that comparison alone whether max-pooling over windows is genuinely weaker for this task, or whether the transformer simply never saw 55\% of its own training set. One reading is an architectural finding; the other is an artefact of our compute budget. They call for very different conclusions, and the difference matters more than any other number in this chapter.

The only way to settle it is to remove the confound, so we retrained the presence transformer on the \textbf{full 53{,}662 window examples} --- everything the training split can produce --- with every other setting identical: same windowing, same negative sampling, same hyperparameters, same seed, same evaluation. The run took 140.7 minutes on the same laptop, against roughly 50 for the subsampled version. Table~\ref{tab:budget} gives the outcome.

\begin{table}[ht]
\centering
\caption{The presence transformer at two training budgets.}
\label{tab:budget}
\small
\begin{tabular}{lccccc}
\toprule
\textbf{Model} & \textbf{Training windows} & \textbf{micro-F1} & \textbf{Precision} & \textbf{Recall}\\
\midrule
Transformer (max-pool)  & 24{,}000 (44.7\%) & 0.6943 & 0.5620 & 0.9080\\
Transformer (max-pool)  & \textbf{53{,}662 (100\%)} & 0.6939 & 0.5546 & 0.9266\\
TF--IDF baseline        & all 408 contracts & 0.7747 & 0.7424 & 0.8101\\
\midrule
\multicolumn{5}{l}{\emph{Paired differences (95\% CI over 10{,}000 contract resamples)}}\\
\multicolumn{2}{l}{Full-data transformer $-$ subsampled transformer} & $-0.0004$ & \multicolumn{2}{l}{$[-0.0098,\ +0.0095]$}\\
\multicolumn{2}{l}{TF--IDF baseline $-$ full-data transformer} & $+0.0809$ & \multicolumn{2}{l}{$[+0.0649,\ +0.0980]$}\\
\bottomrule
\end{tabular}

{\footnotesize\singlespacing \emph{Note.} Evaluated identically on the 102 held-out contracts; differences use the paired cluster bootstrap of Section~\ref{sec:bootstrap}.\par}
\end{table}

\textbf{The confound is not the explanation.} Increasing the training set by 124\% moved micro-F1 by $-0.0004$, with an interval of $[-0.0098, +0.0095]$ squarely straddling zero: the two budgets are statistically indistinguishable. And the gap to the baseline is untouched --- $+0.0809$ on full data against $+0.0805$ on the subsample, with an interval nowhere near zero and the baseline winning in all 10{,}000 resamples. Whatever is holding the windowed transformer back, it is not a shortage of training examples, and the finding this thesis leans on survives the check.

The \emph{shape} of the extra data's effect is worth a sentence, because it is not simply ``nothing happened''. Recall rose from 0.9080 to 0.9266 and precision fell from 0.5620 to 0.5546: more training data made the model fire more readily, not more accurately. That is exactly the failure mode Section~\ref{sec:mil} predicted for max-pooling --- with a mean of 35 windows per contract, the document-level decision is a maximum over many chances to be wrong, and a better-trained window scorer that is still imperfect produces \emph{more} spurious maxima, not fewer. The ceiling here is the aggregation rule, not the encoder underneath it.\footnote{This inference turned out to be only partly right, and Section~\ref{sec:pooling} corrects it: swapping max-pooling for five alternatives moves micro-F1 by at most 0.056 and never past the baseline, so the aggregation rule is a secondary factor rather than the ceiling.}

One secondary effect deserves reporting even though it does not change our conclusion. The AND-ensemble built on the full-data transformer scores 0.7820, and its margin over the baseline is $+0.0073$ with a 95\% interval of $[-0.0000, +0.0140]$ --- it wins in 97.4\% of resamples but the interval still touches zero. So the ensemble's advantage moves from ``clearly a tie'' at the smaller budget to ``borderline'' at the larger one. We are not going to claim it: an interval whose lower bound sits on zero does not clear the bar we set for ourselves in Section~\ref{sec:bootstrap}, and we are not going to lower that bar because the number moved our way. What we will say is that the ensemble's balanced error profile survives the extra data intact.

\textbf{Which model the rest of this chapter uses.} Every other result we report --- the confusion matrix below, the per-category analysis, the end-to-end evaluation of Section~\ref{sec:e2e}, and the deployed Clausify checkpoint --- comes from the 24{,}000-example model. The full-data run is a control experiment, run to test one specific alternative explanation, and we report it as one rather than swapping in the marginally better number after the fact. Since the two are statistically indistinguishable, nothing downstream would change if we had.

\subsection{Ablation: was max-pooling the right aggregation rule?}
\label{sec:pooling}
Section~\ref{sec:mil} chose max-pooling because it \emph{is} the multiple-instance decision rule \cite{maron1998mil} --- a bag is positive if any instance is --- and Sections~\ref{sec:mil} and~\ref{sec:budget} then blamed max-pooling for the transformer's inflated false positives. Both statements were made without ever trying an alternative, which is not good enough for a claim the thesis leans on. We therefore scored every window once and compared aggregation functions on \emph{identical} window scores, so the only thing varying is the rule.

\begin{table}[ht]
\centering
\caption{Aggregation ablation on identical window scores.}
\label{tab:pooling}
\small
\begin{tabular}{lccccc}
\toprule
\textbf{Aggregation} & \textbf{F1@0.5} & \textbf{P@0.5} & \textbf{R@0.5} & \textbf{Best F1} & \textbf{at threshold}\\
\midrule
Max (used in this thesis) & 0.6943 & 0.562 & 0.908 & \textbf{0.7504} & 0.90\\
Top-3 mean                & \textbf{0.7144} & 0.667 & 0.769 & 0.7210 & 0.40\\
Top-5 mean                & 0.6661 & 0.719 & 0.620 & 0.7192 & 0.30\\
Mean over all windows     & 0.1664 & 0.946 & 0.091 & 0.7114 & 0.10\\
Noisy-OR                  & 0.5897 & 0.423 & 0.973 & 0.7022 & 0.95\\
Second-highest window     & 0.6832 & 0.662 & 0.705 & 0.7016 & 0.30\\
\midrule
\emph{TF--IDF baseline (no windows)} & \textbf{0.7747} & 0.742 & 0.810 & --- & ---\\
\bottomrule
\end{tabular}

{\footnotesize\singlespacing \emph{Note.} ``Best'' is each rule's micro-F1 at its own optimal threshold, found by sweeping the test set --- a diagnostic upper bound, not a configuration we adopt.\par}
\end{table}

Three things come out of this, and one of them corrects an earlier claim of ours.

\textbf{At the threshold we actually used, max-pooling was not the best choice.} Top-3 mean scores 0.7144 against max's 0.6943 at the 0.5 boundary --- a fifth of the gap to the baseline, available for free from a one-line change. Averaging the three highest-scoring windows is a natural softening of the MIL rule: it still asks whether evidence is concentrated somewhere in the document, but requires that evidence to appear in more than one place, which is exactly the guard against a single spurious window that Section~\ref{sec:mil} identified as max-pooling's weakness. Mean-over-all-windows collapses at 0.5 (recall 0.091) for the obvious reason --- an average over 35 windows of which one is positive cannot reach 0.5 --- and its best threshold is correspondingly 0.10.

\textbf{Once each rule is given its own threshold, max-pooling is the best of the six.} At 0.90 it reaches 0.7504, ahead of top-3 mean's 0.7210. So max-pooling is not a poor aggregation rule; it is a rule whose output distribution is badly matched to a 0.5 decision boundary, which is the same finding Section~\ref{sec:threshold} reached from the other direction. The two results are one result: the transformer's reported deficit is substantially an operating-point problem, not an aggregation problem.

\textbf{But no aggregation rule reaches the baseline.} The best cell in the entire table --- max-pooling at a threshold tuned on the test set, which is already a generous upper bound --- is 0.7504 against the baseline's 0.7747. Every alternative rule, at every threshold we swept, remains behind. This matters more than the first two points, because it closes off the last available explanation for RQ1: the windowed transformer does not lose because of its training budget (Section~\ref{sec:budget}), or its seed (Section~\ref{sec:seeds}), or its aggregation rule, or its threshold. It loses.

\textbf{A correction to Section~\ref{sec:budget}.} There we wrote that ``the ceiling here is the aggregation rule, not the encoder underneath it'', inferring it from the precision/recall shift when training data was doubled. Table~\ref{tab:pooling} shows that inference was wrong: changing the aggregation rule --- through six variants, each at its best threshold --- moves micro-F1 by at most 0.056 and never past the baseline. The aggregation rule is a real but secondary factor. We have left the original sentence in place with this correction attached, on the same principle as elsewhere in this chapter: an inference we drew from indirect evidence and later measured properly is more useful to a reader than a silently amended one.

A caveat on how far this ablation goes. It varies only the function applied to window scores. It does not vary the window geometry (2{,}000 characters, 1{,}500 stride), which would require re-windowing and retraining, nor does it try a learned aggregator such as attention pooling over window representations. Both remain untested, and attention pooling in particular is the obvious next thing to try, since it could learn to weight windows by relevance instead of taking a fixed statistic.

\subsection{A third source of variance: the split itself}
\label{sec:splitvar}
Sections~\ref{sec:budget} and~\ref{sec:seeds} tested two things that could move our numbers: the training budget and the model seed. Both left the conclusions standing. There is a third source neither touches, and it is arguably the largest: \textbf{the split}. Every result in this thesis comes from one contract-level 80/20 partition drawn with seed 42. With only 510 contracts, a different draw puts different rare categories in the test set and could move the scores on its own. Model-seed variance is not a proxy for this --- reseeding training holds the test set fixed, so it cannot see split effects at all.

Retraining the transformer on twenty splits is not affordable, but the TF--IDF baseline refits in under a minute, and since it is the reference point every comparison in this thesis is made against, its split-to-split spread bounds how much the comparison itself could move. We drew \textbf{20 independent contract-level 80/20 splits} and refitted the baseline on each.

\begin{table}[ht]
\centering
\caption{The TF--IDF baseline refitted on 20 independent contract-level splits.}
\label{tab:splitvar}
\small
\begin{tabular}{lccccc}
\toprule
 & \textbf{Mean} & \textbf{SD} & \textbf{Min} & \textbf{Max} & \textbf{Range}\\
\midrule
micro-F1 over 20 splits & 0.7829 & 0.0063 & 0.7702 & 0.7918 & 0.0216\\
\midrule
\multicolumn{6}{l}{\emph{The split used throughout this thesis (seed 42): 0.7747}}\\
\bottomrule
\end{tabular}
\end{table}

Three readings follow. First, \textbf{split variance is the biggest of the three uncertainties we have measured}: SD 0.0063 against 0.0022 for model seed and a bootstrap interval half-width of about 0.008. It is roughly four and a half times the ensemble's $+0.0014$ margin, which is now the third independent reason to treat that margin as noise.

Second, \textbf{our split is not a lucky one}. At 0.7747 it sits below the 20-split mean of 0.7829, in the lower part of the distribution --- so if anything this thesis reports the baseline at a slightly unfavourable draw. We note this because the alternative would have been worth worrying about: a headline built on a split that happened to flatter the baseline would have made the transformer's defeat look worse than it is.

Third, \textbf{the finding we rely on is far outside this variation}. The baseline's lead over the transformer is 0.0805 --- about thirteen split-level standard deviations. No plausible redraw of the split turns that around.

What we cannot say is anything about \emph{per-category} stability across splits, which is where the instability would really bite: a rare category with two test positives on one draw may have eight on another, and its F1 is meaningless in both cases. Repeated splits or cross-validation for the full pipeline remains future work, and Section~8.2 lists it as a limitation rather than an answered question.

\subsection{What the errors actually look like}
\label{sec:errorexamples}
The confusion matrix counts 311 false positives and 296 false negatives, and counts alone do not tell a reader whether the model is failing stupidly or reasonably. Table~\ref{tab:errors} shows real cases, drawn from the test contracts --- false positives ranked by the model's own confidence, since a system's most confident mistakes are the most revealing ones.

\begin{table}[ht]
\centering
\caption{Representative errors on held-out contracts.}
\label{tab:errors}
\scriptsize
\renewcommand{\arraystretch}{1.25}
\begin{tabular}{@{}p{2.1cm} p{1.5cm} p{9.6cm}@{}}
\toprule
\textbf{Category} & \textbf{Scores} & \textbf{Evidence}\\
\midrule
\multicolumn{3}{@{}l}{\textbf{False positives} --- reported, but the category is not annotated in this contract}\\
\midrule
Insurance & 0.82 / 0.97 & ``\ldots a fixed amount of [*] of gross sales to cover freight, postage, \textbf{insurance costs} on shipments to such Third Party, packing costs\ldots'' \emph{A royalty-deduction clause that mentions insurance as a cost item. The word is present; the obligation is not.}\\
\addlinespace
Governing Law & 0.81 / 0.96 & ``\ldots there shall arise the need to decide the question by what municipal or national law this Agreement \ldots is \textbf{governed}, the following facts shall be excluded\ldots'' \emph{A clause about how to determine the governing law, rather than a clause stating it. Arguably a boundary case in the annotation.}\\
\addlinespace
Parties & 1.00 / 0.99 & An addendum whose opening window reads like a contract header. \emph{Both models fire on document-structural cues; the annotation places Parties elsewhere.}\\
\midrule
\multicolumn{3}{@{}l}{\textbf{False negatives} --- annotated as present, both models scored it low}\\
\midrule
Liquidated Damages & 0.21 / 0.02 & ``\ldots agrees to pay a \textbf{cancellation fee} of Five Thousand Dollars (\$5{,}000) per month multiplied by the number of months remaining\ldots'' \emph{A textbook liquidated-damages term that never uses the phrase. The lexical model has no signal; the transformer, trained on 24 positive examples of this category, has not learned the concept.}\\
\addlinespace
Most Favored Nation & 0.22 / 0.02 & ``\ldots Sponsor shall have the right to approve or disapprove any additional sponsor\ldots unless another\ldots'' \emph{Semantically an MFN-style protection, expressed without any of the category's vocabulary. This category has 21 positive training contracts and scores 0.000 overall.}\\
\addlinespace
Third Party Beneficiary & 0.26 / 0.02 & ``Notwithstanding the foregoing, UTK may assign this Agreement or any portion of its Compensation \ldots to its subsidiaries.'' \emph{This reads like an assignment clause, not a third-party-beneficiary clause. Some of the rare-category ``failures'' are disagreements with the annotation rather than model errors.}\\
\midrule
\multicolumn{3}{@{}l}{\textbf{Span errors} --- category correctly detected, extracted text wrong (Section~\ref{sec:e2e})}\\
\midrule
Covenant Not To Sue & F1 0.00 & Window \emph{does} contain the clause; the model returns ``\texttt{.}'' --- a degenerate single-token span. \emph{A start/end pair collapsing onto punctuation, which the argmax decoding has no constraint against.}\\
\addlinespace
Effective Date & F1 0.06 & Window does \emph{not} contain the gold clause; the model confidently returns ``7th day of september, 1999''. \emph{When Stage 2A supplies the wrong window, Stage 2B answers anyway --- there is no abstention mechanism.}\\
\addlinespace
Insurance & F1 0.20 & Gold is the insurance obligation; the model returns the neighbouring indemnification clause. \emph{Topically adjacent, legally distinct.}\\
\bottomrule
\end{tabular}

{\footnotesize\singlespacing \emph{Note.} Text abridged. Scores are the TF--IDF and transformer probabilities.\par}
\end{table}

Three patterns are worth naming. \textbf{Lexical triggers without the obligation} account for a visible share of false positives --- a contract that discusses insurance costs is not a contract with an insurance covenant, and a full-document bag-of-words model cannot tell the difference. \textbf{Concept-without-vocabulary} explains much of the rare-category failure: a \$5{,}000-per-month cancellation fee \emph{is} liquidated damages, and no model with two dozen training examples of the category will discover that. And \textbf{annotation disagreement} is a real component of the rare tail --- the \emph{Third Party Beneficiary} case above looks like a labelling artefact rather than a miss, which means the 0.000 F1 scores in Section~\ref{sec:raretail} are a mixture of model failure and ceiling effects in the labels. We did not attempt to quantify that mixture; doing so would need a re-annotation exercise we are not qualified to run.

The span errors point somewhere more actionable. The degenerate ``\texttt{.}'' prediction and the confident answer from a window that cannot contain the answer are both symptoms of the same missing piece: the extractor has no way to say \emph{nothing here}. SQuAD~v2-style unanswerable training, which supplies exactly that, is the concrete fix and it is listed in future work.

\subsection{Confusion-matrix analysis}
Our evaluation protocol called for a full confusion-matrix analysis of the final model. Table~\ref{tab:cm} gives the matrix for the AND-ensemble over all test decisions, Figure~\ref{fig:cm} shows the same matrix as a heatmap, and Table~\ref{tab:metrics} lists every metric derived from it.

\begin{table}[ht]
\centering
\caption{Confusion matrix of the AND-ensemble on the 20\% test set (4{,}182 decisions).}
\label{tab:cm}
\begin{tabular}{lcc}
\toprule
 & \textbf{Predicted Absent} & \textbf{Predicted Present}\\
\midrule
\textbf{Actually Absent}  & TN = 2{,}523 & FP = 311\\
\textbf{Actually Present} & FN = 296     & TP = 1{,}052\\
\bottomrule
\end{tabular}
\end{table}

\begin{table}[ht]
\centering
\caption{Test metrics derived from the confusion matrix.}
\label{tab:metrics}
\begin{tabular}{llc}
\toprule
\textbf{Metric} & \textbf{Formula} & \textbf{Value}\\
\midrule
Accuracy    & $(TP+TN)/\text{all}$          & 85.49\%\\
Precision   & $TP/(TP+FP)$                  & 77.18\%\\
Recall (Sensitivity) & $TP/(TP+FN)$         & 78.04\%\\
Specificity & $TN/(TN+FP)$                  & 89.03\%\\
F1 score    & $2PR/(P+R)$                   & 0.7761\\
\bottomrule
\end{tabular}
\end{table}

\begin{figure}[ht]
\centering
\includegraphics[width=0.6\textwidth]{confusion_matrix.png}
\caption{Confusion matrix of the AND-ensemble on the test set.}
\label{fig:cm}
\end{figure}

Precision and recall come out balanced (77.2\% and 78.0\%), which tells us the AND rule did its job: raw max-pooling gives up precision, and requiring agreement wins it back without sacrificing recall. The specificity of 89.0\% means the system rarely claims a clause that is not there --- in a legal-review assistant, that is the more damaging kind of mistake.

\subsection{Per-category behaviour}
Figure~\ref{fig:percat} shows per-category F1 for the strongest categories. The frequent, formulaic ones are close to solved (\emph{Document Name} 1.000, \emph{Parties} 0.995, \emph{Governing Law} 0.961, \emph{Agreement Date} 0.958). At the other end, the five categories with at most seven test positives (\emph{Most Favored Nation}, \emph{No-Solicit of Customers}, \emph{Price Restrictions}, \emph{Source Code Escrow}, \emph{Unlimited License}) all score 0.000. That number needs context: with two positive examples in the test set, a score of zero means the model missed two contracts, not that it is broken. This tail is a property of the dataset rather than of any architecture, and it is why we quote micro- rather than macro-averaged numbers in our headline results.

\begin{figure}[ht]
\centering
\includegraphics[width=0.85\textwidth]{per_category_f1.png}
\caption{Per-category F1 for the twelve strongest categories.}
\label{fig:percat}
\end{figure}

\subsection{Does the AND rule punish rare categories harder than common ones?}
\label{sec:raretail}
Section~\ref{sec:bootstrap} defended the conjunctive ensemble on its error profile: it converts the transformer's inflated false positives into balanced precision and recall. That defence was made on the aggregate, and an aggregate can hide a trade that falls unevenly. It is worth asking where the AND rule's precision comes from, because there is an obvious reason to suspect the rare tail specifically. Requiring two models to agree is a stricter bar exactly where each model individually has the least training signal to be confident with --- so the categories with the fewest positive examples are the ones where the second model is least likely to concur, and agreement is what the rule demands.

We split the 41 categories by \emph{training} frequency, took the bottom ten (11 to 42 positive training contracts, 70 positive test decisions between them), and compared the AND rule against the transformer alone on that tail and on the remaining 31 categories separately. Table~\ref{tab:raretail} gives the result, and it is not the same story in the two halves.

\begin{table}[ht]
\centering
\caption{The ensemble on the ten rarest categories versus the rest.}
\label{tab:raretail}
\small
\begin{tabular}{llccccc}
\toprule
\textbf{Group} & \textbf{Model} & \textbf{Precision} & \textbf{Recall} & \textbf{F1} & \textbf{TP} & \textbf{FN}\\
\midrule
\multirow{3}{*}{Rarest 10 (70 test positives)}
 & Transformer & 0.261 & 0.500 & 0.343 & 35 & 35\\
 & TF--IDF     & 0.429 & 0.300 & 0.353 & 21 & 49\\
 & \textbf{AND ensemble} & 0.586 & 0.243 & 0.343 & 17 & 53\\
\addlinespace
\multirow{3}{*}{Other 31 (1{,}278 test positives)}
 & Transformer & 0.582 & 0.930 & 0.716 & 1{,}189 & 89\\
 & TF--IDF     & 0.753 & 0.838 & 0.793 & 1{,}071 & 207\\
 & \textbf{AND ensemble} & 0.776 & 0.810 & \textbf{0.792} & 1{,}035 & 243\\
\bottomrule
\end{tabular}

{\footnotesize\singlespacing \emph{Note.} Micro-averaged within each group.\par}
\end{table}

\textbf{On the common categories the ensemble works as advertised.} Against the transformer alone it gains 19.4 points of precision for 12.1 points of recall, and F1 rises from 0.716 to 0.792. That is a good trade, and it is where the ensemble's entire aggregate advantage is generated.

\textbf{On the rare tail it is not a free lunch --- it is a straight swap at no gain.} The AND rule buys 32.5 points of precision by giving up 25.7 points of recall, and the F1 lands at 0.343: identical, to three decimals, to the transformer alone. Every point of precision on the rare tail is paid for with recall, and nothing is left over. In absolute terms the transformer alone finds 35 of the 70 rare positives; requiring agreement cuts that to 17. Notably, the plain TF--IDF baseline is the best of the three on this tail (F1 0.353), which is a further reminder of how little the transformer contributes where training signal is thin.

The per-category detail is starker than the micro-average. On five of the ten rarest categories the AND rule scores exactly 0.000, and on three of those --- \emph{Price Restrictions}, \emph{Unlimited/All-You-Can-Eat License}, \emph{No-Solicit of Customers} --- the transformer alone was finding something (F1 0.095, 0.667, 0.286) before the conjunction zeroed it out. The rule does not merely fail to help on the rare tail; on those categories it removes detections that one component had made.

\textbf{Why this matters for deployment rather than just for the table.} The aggregate micro-F1 weights every one of the 4{,}182 decisions equally, but a contract signer does not. Missing a rare clause is not obviously less costly than missing a common one, and it may be considerably more so: \emph{Most Favored Nation} sits in the rarest ten and is one of our eight High-risk categories, and both models score 0.000 on it. A system whose ensemble rule is tuned --- even by a parameter-free rule --- toward precision on frequent categories is optimizing the metric rather than the harm. We report the ensemble as our headline configuration because it is what the aggregate protocol selects, but on this evidence a deployment aimed at protecting a signer would be better served by \emph{running the more permissive rule on the High-risk categories and the conjunction elsewhere}, accepting more false alarms exactly where a missed clause costs the most. We have not implemented that split rule --- doing so and evaluating it properly is future work (Section~8.3), not something to bolt on after seeing the test set.


\section{The Threshold and What the Confidence Score Means}
\label{sec:threshold}

Every presence decision in this thesis, and every clause card in Clausify, comes from comparing a probability against 0.5. Section~\ref{sec:valprotocol} records that this threshold was never chosen --- it is the default decision boundary of a two-class softmax, inherited rather than selected. Two questions follow, and neither had been asked: is 0.5 a sensible operating point, and does the number being compared against it mean anything? Figure~\ref{fig:threshcal} answers both.

\begin{figure}[ht]
\centering
\includegraphics[width=\textwidth]{threshold_calibration.png}
\caption{Threshold sweep and reliability diagram.}
\label{fig:threshcal}
\end{figure}

\subsection{0.5 is a poor operating point for the transformer}
Sweeping the threshold post hoc shows the max-pooled transformer reaches micro-F1 \textbf{0.750 at a threshold of 0.9}, against \textbf{0.694} at the default --- a gain of 0.056, which is most of the way to closing its gap with the TF--IDF baseline (0.775). The reason is the one Section~\ref{sec:mil} predicted structurally: a maximum over ~35 windows is inflated by construction, so the scores pile up near 1 and a boundary drawn at 0.5 admits far too much. Precision at 0.5 is 0.562; at 0.9 it is 0.728.

\textbf{We are not adopting 0.9.} It was found by looking at the test set, and adopting it would convert a held-out evaluation into a tuned one --- the precise failure Section~\ref{sec:valprotocol} certifies this project against. The honest statement is narrower and more useful: \emph{a meaningful part of the transformer's reported deficit is an un-chosen threshold rather than an inability to discriminate}, and our comparison in Section~\ref{sec:bootstrap} therefore compares a tuned-by-construction baseline against an untuned transformer. A fair rematch selects the threshold for each model on a validation split carved from the 408 training contracts, and we did not do that. Section~\ref{sec:budget} showed the gap is not a data-budget artefact; this section shows part of it may be a threshold artefact, and that qualification belongs on the finding.

The AND ensemble behaves oppositely: 0.5 is already its optimum (0.776), and raising the threshold collapses it --- at 0.7 micro-F1 falls to 0.610 and High-risk recall to 0.159, because requiring two models to \emph{both} clear a high bar is close to requiring unanimity among the confident. The dashed curves in Figure~\ref{fig:threshcal} make the tension visible: for the transformer, the threshold that maximizes micro-F1 (0.9) drives High-risk recall down from 0.841 to 0.688. There is no single threshold that is simultaneously best for the aggregate metric and for the clauses that matter most.

\subsection{The confidence bar in the application is not a probability}
Clausify displays the presence score to the user as a confidence bar. The reliability diagram tests whether that display means what a user would take it to mean, and it does not.

The max-pooled transformer has an \textbf{Expected Calibration Error of 0.201} and a Brier score of 0.194. It is overconfident across the entire range: among predictions in the $[0.5, 0.6)$ bin the clause is actually present \textbf{19.5\%} of the time, and even in the top bin --- scores above 0.9, which the interface renders as a nearly full bar --- the clause is present \textbf{72.8\%} of the time. A user reading ``96\% confident'' is looking at something closer to a 73\% chance. Again the cause is max-pooling: taking the maximum over dozens of windows is not a probability-preserving operation, and nothing in training asked it to be.

The TF--IDF baseline is differently miscalibrated (ECE 0.172) but in the safer direction over most of its range --- it is \emph{under}confident in the middle, with the $[0.7, 0.8)$ bin correct 89.6\% of the time and the $[0.8, 0.9)$ bin 97.8\%.

The practical consequence for the deployed system is direct and we have recorded it in Section~\ref{sec:demolimits}: the threshold slider, which the application offers as the user's control over conservativeness, is calibrated in units that do not correspond to probability of correctness. Fixing this needs no retraining --- Platt scaling or isotonic regression fitted on a validation split would map scores onto probabilities --- and it is the cheapest single improvement available to the interface. We did not have a validation split set aside to fit it on, which is a consequence of the protocol gap in Section~\ref{sec:valprotocol} rather than an independent oversight.

\section{Risk-Weighted Evaluation: Does the Metric Match the Purpose?}
\label{sec:riskeval}

Micro-F1 over 4{,}182 decisions treats every decision as equally important. A contract signer does not. The whole premise of the risk layer (Section~\ref{sec:risk}) is that some categories matter more than others --- so the evaluation should ask whether the system's errors fall where they do least harm. We had not asked, and the answer turns out to reverse the recommendation the aggregate produces.

Table~\ref{tab:riskeval} splits the 41 categories into the three risk bands of our own taxonomy and scores each model within each band.

\begin{table}[ht]
\centering
\caption{Presence performance within each risk band.}
\label{tab:riskeval}
\small
\begin{tabular}{llccccc}
\toprule
\textbf{Band} & \textbf{Model} & \textbf{Precision} & \textbf{Recall} & \textbf{F1} & \textbf{Found} & \textbf{Missed}\\
\midrule
\multirow{4}{*}{\shortstack[l]{\textbf{High}\\8 cats\\176 positives}}
 & Transformer  & 0.404 & \textbf{0.841} & 0.546 & 148 & \textbf{28}\\
 & TF--IDF      & 0.540 & 0.648 & 0.589 & 114 & 62\\
 & \textbf{AND ensemble} & 0.585 & 0.625 & 0.604 & 110 & \textbf{66}\\
 & OR ensemble  & 0.391 & \textbf{0.864} & 0.538 & 152 & \textbf{24}\\
\addlinespace
\multirow{4}{*}{\shortstack[l]{\textbf{Medium}\\22 cats\\498 positives}}
 & Transformer  & 0.425 & 0.871 & 0.571 & 434 & 64\\
 & TF--IDF      & 0.621 & 0.723 & 0.668 & 360 & 138\\
 & \textbf{AND ensemble} & 0.647 & 0.673 & 0.659 & 335 & 163\\
 & OR ensemble  & 0.423 & 0.922 & 0.580 & 459 & 39\\
\addlinespace
\multirow{4}{*}{\shortstack[l]{\textbf{Low}\\11 cats\\674 positives}}
 & Transformer  & 0.813 & 0.953 & 0.877 & 642 & 32\\
 & TF--IDF      & 0.909 & 0.917 & 0.913 & 618 & 56\\
 & \textbf{AND ensemble} & 0.924 & 0.901 & 0.912 & 607 & 67\\
 & OR ensemble  & 0.803 & 0.969 & 0.878 & 653 & 21\\
\bottomrule
\end{tabular}

{\footnotesize\singlespacing \emph{Note.} Threshold 0.5. ``Missed'' counts gold clauses of that band the model failed to report.\par}
\end{table}

\subsection{The headline configuration is the worst one for High-risk clauses}
Read down the High band. The AND ensemble --- the configuration this thesis reports as its headline, on the strength of aggregate micro-F1 --- \textbf{misses 37.5\% of High-risk clauses}. The transformer alone, the component we have spent two sections establishing is the weaker model, misses \textbf{15.9\%}. The permissive OR rule, which finishes \emph{last} on aggregate micro-F1 (0.696), misses only \textbf{13.6\%}.

In absolute terms, on 176 High-risk gold clauses in 102 contracts: the ensemble surfaces 110 and loses 66; OR surfaces 152 and loses 24. Choosing the ensemble over OR on the basis of micro-F1 costs a signer 42 High-risk clauses --- uncapped liabilities, IP assignments, non-competes --- in exchange for a leaderboard position that Section~\ref{sec:bootstrap} already showed is a statistical tie.

The mechanism is the one identified in Section~\ref{sec:raretail} and it compounds here. High-risk categories are disproportionately rare (\emph{Most Favored Nation}, \emph{Irrevocable/Perpetual License}, \emph{Minimum Commitment} all sit in the sparse tail), the conjunction is hardest to satisfy exactly where training signal is thinnest, and the risk band that most needs recall is the one where requiring agreement destroys it. The Low band shows the inverse: 674 of the 1{,}348 gold clauses are Low-risk --- dates, party names, governing law --- and every model scores near 0.9 there. \textbf{Aggregate micro-F1 is dominated by the clauses that matter least.}

\subsection{What we conclude, and what we do not}
We do not switch the reported configuration on the strength of this table, and it is worth being explicit about why: choosing OR \emph{after} seeing which rule scores best on the High band of the test set is test-set selection, exactly the practice Section~\ref{sec:valprotocol} says this project avoided. The result stands as a diagnosis, not a retuning.

What it does establish is that \textbf{our evaluation metric and our stated purpose were never aligned}, and that the misalignment is large enough to invert the model ranking. A system whose declared aim is to warn a signer before they sign should be selected on High-risk recall, or on a risk-weighted F1, not on an average over 41 categories in which two thirds of the mass is dates and party names. Section~8.3 lists the proper version of this experiment: define the risk-weighted objective first, then select the rule on a validation split, then report on the test set once.

Finally, this is a further argument for taking Section~\ref{sec:taxvalid} seriously. Every number in Table~\ref{tab:riskeval} is computed against \emph{our own} assignment of categories to bands. If a lawyer disagrees about which categories are High, the table changes with it. The risk-weighted evaluation inherits all the unvalidated-taxonomy caveats and adds no independent evidence to them.

\section{Span Extraction}
\label{sec:span}
We evaluated the QA model on all 2{,}667 answer-bearing test windows (Table~\ref{tab:span}). A token-F1 of 0.763, with 82.7\% of predictions overlapping the gold span at F1 $\geq$ 0.5, means that when the clause is in front of the model, it captures most of the right text roughly eight times out of ten. The lower exact-match score did not surprise us: annotated legal spans are long, and their exact boundaries are genuinely debatable (should a leading section number be included?), so content overlap is the metric that carries information here. For scale, the published CUAD baselines at this task use DeBERTa-v3-large, roughly five times the parameters of our distilled backbone.

Section~\ref{sec:bootstrap} argued that a point estimate on 102 contracts should not be read without an interval, and that argument does not stop applying when the task changes. We therefore ran the \emph{same} cluster bootstrap here. The correlation structure is if anything stronger than in presence classification: the 2{,}667 span decisions come from only 102 contracts, and a contract with unusual drafting conventions --- numbered sub-clauses, defined-term-heavy prose --- shifts every span in it in the same direction. Table~\ref{tab:span} therefore reports two intervals per metric: one resampling \emph{contracts} (the correct unit) and one resampling \emph{clauses} independently (the naive unit), because the gap between them is itself informative.

\begin{table}[ht]
\centering
\caption{Span extraction on the answer-bearing test windows.}
\label{tab:span}
\small
\begin{tabular}{lccc}
\toprule
\textbf{Metric} & \textbf{Score} & \textbf{95\% CI (by contract)} & \textbf{95\% CI (by clause)}\\
\midrule
Token-level F1                & 0.763 & $[0.740,\ 0.782]$ & $[0.751,\ 0.774]$\\
Overlap (token-F1 $\geq$ 0.5) & 82.7\% & $[79.7\%,\ 85.1\%]$ & $[81.3\%,\ 84.1\%]$\\
Exact Match                   & 35.5\% & $[32.5\%,\ 38.6\%]$ & $[33.7\%,\ 37.3\%]$\\
\bottomrule
\end{tabular}

{\footnotesize\singlespacing \emph{Note.} 95\% bootstrap intervals, 10{,}000 resamples. The contract-clustered interval is the one to quote; the clause-level interval is shown only to expose how far a naive resampling unit understates the uncertainty.\par}
\end{table}

The contract-clustered interval on token-F1 is $[0.740, 0.782]$ --- roughly \emph{twice} as wide as the $[0.751, 0.774]$ that treating clauses as independent would have produced. That factor is the practical cost of ignoring clustering, and it is the same effect that made the cluster bootstrap necessary in Section~\ref{sec:bootstrap}. It also sets the resolution of this number: differences below about two points of token-F1 are not measurable on a test set of this size, which is worth remembering before comparing 0.763 against any published figure.

\textbf{Scope caveat.} These numbers measure extraction quality \emph{given an answer-bearing window}. The full-document benchmark in the CUAD paper (precision at 80\% recall over all windows, including the answer-free ones) needs a heavier document-stride ranking evaluation that we did not run, so we make no claim on that metric. Section~\ref{sec:e2e} measures the complementary thing: what happens to span quality when the window is chosen by the presence stage instead of being handed over.

\section{Summarization}
\label{sec:summ}

\subsection{What the reference set actually is}
Before any number in this section can be read correctly, one property of the experiment has to be stated plainly, because it determines what the metric can possibly mean. CUAD contains no plain-English rewrites, so there was no annotated target set to train or test against. We therefore \textbf{wrote the targets ourselves}: 41 single-sentence summaries, one per category, phrased from the weaker party's perspective (Section~4.4). Every training pair and every evaluation pair reuses one of those 41 sentences. There is no clause-level variation in the target set at all --- two \emph{Cap on Liability} clauses with completely different caps share a target string, character for character.

That is a design choice with a consequence we should own rather than discover: with 41 unique targets covering 2{,}667 test clauses, a model can score well by learning to \emph{recognise} the category and reproduce its sentence, and the metric cannot tell that apart from summarization. The experiment below is built to separate the two.

\subsection{Two reference sets, the same 150 clauses}
We evaluate the fine-tuned FLAN-T5-small on 150 held-out test clauses (sampled with seed 42) against two different reference sets:

\begin{enumerate}[itemsep=2pt]
  \item \textbf{Template references} --- the category template, exactly as in training. This is the protocol behind the headline number.
  \item \textbf{Clause-specific references} --- 150 one-sentence summaries, one per clause, written to describe \emph{that} clause: its actual cap, its actual notice period, its actual parties. These were drafted with a large-language-model assistant working from the clause text alone, without sight of any model output, and are deliberately independent of the 41 templates. They are a pilot reference set, not lawyer-authored annotation, and we make no claim beyond that.
\end{enumerate}

Because the clauses are identical across the two conditions, the difference between the two scores isolates exactly one thing: how much of the reported ROUGE-L comes from reference-set uniformity rather than from summarization. Table~\ref{tab:sumeval} gives the result.

\begin{table}[ht]
\centering
\caption{ROUGE-L under two reference sets, same 150 clauses.}
\label{tab:sumeval}
\small
\begin{tabular}{llcc}
\toprule
\textbf{System} & \textbf{Reference set} & \textbf{ROUGE-L} & \textbf{95\% CI}\\
\midrule
Fine-tuned FLAN-T5-small & 41 category templates    & 0.775 & $[0.718,\ 0.833]$\\
Fine-tuned FLAN-T5-small & clause-specific (pilot)  & 0.116 & $[0.105,\ 0.128]$\\
Zero-shot FLAN-T5-small  & clause-specific (pilot)  & 0.299 & $[0.275,\ 0.322]$\\
\midrule
\emph{Category template itself} & clause-specific (pilot) & 0.115 & $[0.104,\ 0.127]$\\
\bottomrule
\end{tabular}

{\footnotesize\singlespacing \emph{Note.} Intervals are 95\% bootstrap CIs over the 150 clauses (10{,}000 resamples).\par}
\end{table}

\subsection{Reading the table}
Three things follow, and none of them is a hedge.

\textbf{The 0.775 is template reproduction, and we can now say so quantitatively.} On all 150 test clauses the fine-tuned model emits one of the 41 template sentences \emph{verbatim} --- 100\% of outputs are an exact copy of a template, the correct one 70\% of the time and a different category's template the other 30\%. The model never produces a sentence of its own. When it is wrong it is wrong in the way a classifier is wrong: a \emph{Price Restrictions} clause about a 10\% rate cap is answered with the \emph{ROFR} template. Fine-tuning did not teach the model to summarize; it taught it to pick one of 41 strings.

\textbf{The last row of Table~\ref{tab:sumeval} closes the argument.} The bare category template, with no model involved at all, scores 0.115 against the clause-specific references --- statistically indistinguishable from the fine-tuned model's 0.116. A constant lookup table achieves the model's entire score. Whatever the fine-tuned summarizer has learned adds nothing measurable on top of ``print this category's sentence''.

\textbf{Fine-tuning made the model worse at the actual task.} Against clause-specific references the \emph{zero-shot} model scores 0.299, more than twice the fine-tuned model's 0.116. We do not present this as evidence that zero-shot FLAN-T5 is a good legal summarizer --- it is not; as Table~\ref{tab:sumqual} shows, it largely echoes the source text, and echoing inflates a longest-common-subsequence metric against a reference that necessarily shares the clause's vocabulary. The honest reading is narrower and still uncomfortable: our fine-tuning traded away whatever clause-specific content the base model had, in exchange for fluent category-level boilerplate that scores well only against a reference set built from the same boilerplate.

\begin{table}[ht]
\centering
\caption{Fine-tuned vs.\ zero-shot summarizer on the same test clause (abridged).}
\label{tab:sumqual}
\small
\begin{tabular}{p{3.2cm} p{10.3cm}}
\toprule
\textbf{Input clause} & ``Any attempted assignment or delegation in violation of this section by either party without the prior written consent of the other will be void.''\\
\midrule
Gold template & You cannot transfer this contract to someone else without permission.\\
Fine-tuned    & You cannot transfer this contract to someone else without permission. \emph{(exact template)}\\
Zero-shot     & ``a) Any attempted assignment or delegation in violation of this section \ldots will be void.'' \emph{(echoes the legalese)}\\
\midrule
Clause-specific reference & Neither side may hand its rights or duties to a third party without the other's written consent, which cannot be unreasonably refused.\\
\bottomrule
\end{tabular}
\end{table}

\subsection{What we therefore claim}
We claim a working \textbf{classification-to-template} system: it identifies a clause's category and attaches a pre-written, human-authored plain-English sentence for that category, which is genuinely useful in the Clausify interface --- a signer reading ``You cannot transfer this contract to someone else without permission'' next to a wall of legalese is better off than before. We do \textbf{not} claim free-text legal summarization, and the 0.775 should never be quoted as evidence of it. The number that describes clause-level summarization ability is 0.116, and the honest conclusion is that this component is a proof of concept, not a result.

Getting past this needs a target set with clause-level variation, which is exactly what the pilot above provides a small taste of: 150 references took a few hours to draft. Section~8.3 sets out the scaled version.

\section{End-to-End: What Actually Reaches the User}
\label{sec:e2e}

Every number so far has been measured with the other stages held out of the way. Presence classification was scored given the ground-truth category list; span extraction was scored given a window already known to contain the answer; summarization was scored given a gold clause. Each of those is the right way to evaluate a component in isolation, and none of them describes what Clausify does to a contract. The stages run in series, and errors compound: a category the presence stage misses is a clause the span stage never sees.

So we ran the pipeline end to end on all 102 test contracts and asked one question of each of the 1{,}348 gold clauses: does it survive \emph{all three} stages? A clause counts as surviving only if (i) the presence stage detects its category, (ii) the span stage --- reading the window the presence stage selected, never a window chosen using the gold offset --- returns text overlapping the gold clause at token-F1 $\geq$ 0.5, and (iii) the risk level attached is the correct one. Table~\ref{tab:e2e} gives the attrition.

\begin{table}[ht]
\centering
\caption{End-to-end survival of the 1{,}348 gold clauses.}
\label{tab:e2e}
\small
\begin{tabular}{llcc}
\toprule
\textbf{Stage} & \textbf{Condition} & \textbf{Clauses surviving} & \textbf{Share of gold}\\
\midrule
--- & gold clauses in the test set & 1{,}348 & 100.0\%\\
2A Presence & category detected & 1{,}052 & 78.0\%\\
2B Span & extracted text overlaps gold at F1 $\geq$ 0.5 & 293 & 21.7\%\\
Risk & correct level attached & 293 & 21.7\%\\
\midrule
\multicolumn{2}{l}{\textbf{End-to-end survival}} & \textbf{293} & \textbf{21.7\%}\\
\multicolumn{2}{l}{\emph{95\% CI, cluster bootstrap over contracts}} & & $[19.6\%,\ 23.9\%]$\\
\bottomrule
\end{tabular}

{\footnotesize\singlespacing \emph{Note.} AND-ensemble presence rule. The last column is the fraction of clauses still alive after that stage.\par}
\end{table}

\textbf{21.7\% is the number that describes a user's experience}, and it is far below what the component scores would lead a reader to expect. Multiplying the isolated figures --- 78.0\% presence recall times 82.7\% span overlap --- predicts about 64\%. The pipeline delivers a third of that. The gap is the whole point of running this evaluation, so it deserves to be explained rather than reported.

The risk stage is the one place where nothing is lost: the level is a fixed per-category lookup, so it is correct exactly when the category is, and it introduces no error of its own. That is a property of the design, not an achievement --- and it is worth remembering that ``correct'' here means ``matches our taxonomy'', which Section~\ref{sec:taxvalid} notes has not been externally validated.

\subsection{Where the clauses are lost}
The collapse happens at stage 2B, and it has two distinct causes, which we separated by checking, for every detected clause, whether the window the presence model selected actually contained the gold text.

\begin{table}[ht]
\centering
\caption{Decomposing the span-stage loss over the 1{,}052 clauses whose category was detected.}
\label{tab:e2eattr}
\small
\begin{tabular}{lcc}
\toprule
 & \textbf{Count} & \textbf{Share}\\
\midrule
Selected window \emph{contains} the gold clause & 718 & 68.3\%\\
\quad --- and the span model then found it (F1 $\geq$ 0.5) & 285 & 39.7\% of 718\\
\quad --- and the span model missed it anyway & 433 & 60.3\% of 718\\
Selected window does \emph{not} contain the gold clause & 334 & 31.7\%\\
\quad --- span scored $\geq$ 0.5 regardless (duplicate wording elsewhere) & 8 & \\
\bottomrule
\end{tabular}
\end{table}

\textbf{Cause one: the max-pooled window is often the wrong window.} Max-pooling was chosen (Section~\ref{sec:mil}) because it removes the recall ceiling --- if any window fires, the category is reported --- and Section~\ref{sec:mil} was careful about what that costs in \emph{precision}. What it did not anticipate is that the argmax window is also being used as a \emph{locator}, and it is a poor one: 31.7\% of the time the highest-scoring window does not contain the clause at all. The MIL objective never asked it to. The model was trained to answer ``does this window contain a clause of this category'', and a window that discusses liability in general terms can outscore the one that actually caps it. For those 334 clauses the span stage is being handed the wrong 2{,}000 characters, and no extractor could recover them.

\textbf{Cause two: the span model is much weaker on windows it was not trained to see.} Even with a window that does contain the clause, extraction succeeds only 39.7\% of the time, against 82.7\% in the isolated benchmark of Section~\ref{sec:span}. The reason is the focus-window construction of Section~4.3. At training time every example was re-centred so the answer began roughly 150 characters into a 1{,}200-character window. That is legitimate --- it uses a label to build a training example, never an inference input, and Section~\ref{sec:leakage} defends it on exactly those grounds. But legitimacy is not the same as harmlessness: the model has only ever read windows in which the answer sits near the start, and at inference it gets blind 1{,}200-character chunks in which the answer may be anywhere. Mean token-F1 falls from 0.763 to 0.386. The re-centring did not leak the test set; it did quietly narrow the distribution the model can handle.

\subsection{What this changes}
Three things follow, and we would rather state them than let the component tables stand unqualified.

First, \textbf{the span numbers in Section~\ref{sec:span} are an upper bound on deployed extraction quality, not a description of it}. They answer ``how good is the extractor when handed the right window'', which is a fair question about the model and a misleading one about the system.

Second, \textbf{the most valuable next engineering step is not a bigger backbone}. Cause two is a data-construction problem with a cheap fix --- train the span model on windows where the answer is placed at varied offsets, or on the same blind chunks it will see at inference --- and cause one argues for choosing the extraction window by a locator trained for that job rather than reusing a presence score. Both are within the compute budget of this project; neither requires more parameters.

Third, \textbf{Clausify's output should be read as a shortlist, not an extraction}. When the application shows a clause card, the category is right about 78\% of the time it was present, but the quoted text is a good rendering of the actual clause in roughly one case in five overall. The application already tells users it is not legal advice; on this evidence it should also tell them that the quoted text may not be the clause the badge refers to, and Section~\ref{sec:demolimits} now says so.

\section{Summary Scorecard}
Table~\ref{tab:scorecard} pulls together the held-out results of all three models. Every number was produced by the standalone test notebook running on the saved 80/20-trained models.

\begin{table}[ht]
\centering
\caption{Final held-out test results (408 training / 102 test contracts).}
\label{tab:scorecard}
\small
\begin{tabular}{llll}
\toprule
\textbf{Stage} & \textbf{Model} & \textbf{Metric} & \textbf{Score}\\
\midrule
2A Presence & AND-ensemble (TF--IDF + DistilBERT) & micro-F1  & 0.776\\
2A Presence & AND-ensemble & Accuracy               & 85.49\%\\
2A Presence & AND-ensemble & Precision / Recall     & 77.18\% / 78.04\%\\
2A Presence & AND-ensemble & Specificity            & 89.03\%\\
2B Span     & DistilBERT-QA & token-F1 (given the window) & 0.763\\
2B Span     & DistilBERT-QA & token-F1 (window chosen by 2A) & 0.386\\
2B Span     & DistilBERT-QA & Exact Match           & 35.5\%\\
1 Summarizer & FLAN-T5-small & ROUGE-L (41 templates)  & 0.775\\
1 Summarizer & FLAN-T5-small & ROUGE-L (clause-specific) & 0.116\\
\midrule
\textbf{Full pipeline} & \textbf{all three stages} & \textbf{gold clauses surviving} & \textbf{21.7\%}\\
\bottomrule
\end{tabular}
\end{table}

\section{Discussion: What the Evidence Supports}
Pulling the chapter together, the results fall into three groups, and it matters which is which.

\textbf{Findings that survived every check we could run.}
\begin{enumerate}[itemsep=2pt]
  \item For document-level \emph{presence}, a full-document lexical baseline beats a windowed DistilBERT by $+0.081$ micro-F1. We tested the two obvious alternative explanations and neither held: training on all 53{,}662 windows instead of 24{,}000 changed the score by $-0.0004$ (Section~\ref{sec:budget}), and across three seeds the gap is 9.5 standard deviations wide (Section~\ref{sec:seeds}). More data made the transformer fire \emph{more} rather than better, which is what max-pooling over 35 windows predicts. The limitation is the aggregation rule, not the encoder or the budget.
  \item For \emph{span extraction}, the transformer has no substitute --- a bag-of-words model cannot point at character offsets --- and it performs well in distilled form \emph{when handed the right window}: token-F1 0.763, 95\% CI $[0.740, 0.782]$.
  \item For \emph{summarization}, a 41-sentence template seed turns the task into classification. The fine-tuned model reproduces a template verbatim on 100\% of test clauses, and against clause-specific references it scores no better than the template alone (Section~\ref{sec:summ}). A high automatic metric should not be believed until the outputs have been read.
\end{enumerate}

\textbf{A claim we withdrew.} The AND-ensemble's first place in Table~\ref{tab:ensemble} is not a result. Its margin over the baseline is $+0.0014$; the cluster bootstrap puts that inside $[-0.0068, +0.0089]$ (Section~\ref{sec:bootstrap}), and the seed-to-seed standard deviation of the ensemble itself is $0.0022$ --- larger than the margin, with a range three times larger. Two independent sources of uncertainty agree that the top two rows of that table are a tie. What the ensemble does earn is a balanced error profile and a fourfold reduction in seed variance relative to the transformer alone.

\textbf{Costs we did not see until we looked for them.} Two results in this chapter came from questions we had not thought to ask, and both are unflattering.
\begin{itemize}[itemsep=2pt]
  \item The conjunctive rule is not uniformly good. On the ten rarest categories it nets \emph{zero} F1 gain over the transformer alone, buying 32.5 points of precision with 25.7 points of recall and cutting detections from 35 of 70 to 17 of 70 (Section~\ref{sec:raretail}). Its entire aggregate benefit is generated on common categories, and one of the categories it zeroes out is High-risk.
  \item The stage-wise scores do not compose. Only 21.7\% of gold clauses survive the full pipeline, against roughly 64\% predicted by multiplying the component numbers (Section~\ref{sec:e2e}), because the presence model's argmax window is the wrong window 31.7\% of the time and the span model was trained on re-centred windows it never sees at inference.
\end{itemize}

The through-line is methodological rather than architectural. Every number in this chapter that we originally trusted least --- the ensemble's rank, the summarizer's ROUGE, the span model's token-F1 --- turned out to be measuring something narrower than it appeared to, and in each case the way to find that out was to measure the thing the metric was standing in for: resample the contracts, change the seed, vary the reference set, run the stages in series. None of those checks needed a bigger model or a larger dataset. They needed the question to be asked.
